\section{Conceptual framework}
\label{sec:framework}

\subsection{Three layers, three questions}

The methodology is organised into three layers, each answering a distinct
question that a planner asks in sequence. The structure is shown in
\Cref{fig:framework}.

\begin{description}[leftmargin=0pt, labelwidth=0pt, labelsep=0pt, style=nextline]
  \item[Layer 1 --- Baseline assessment: \emph{does this site deserve
  attention?}] A Heat Exposure Score and a Vulnerability Score are computed
  from the site description and combined into a Heat Priority Index. This layer
  is independent of the proposed intervention: it characterises the place, not
  the project. Two options assessed for the same site share an identical
  Layer 1 result, which is what makes side-by-side comparison meaningful.

  \item[Layer 2 --- NbS performance: \emph{what can this intervention deliver
  here?}] A base value drawn from the typology library is modulated by a site
  adjustment factor derived from conditions the user has already described.
  This layer produces the Cooling Potential Score and the indicative
  temperature reduction range. It is where the tool's evidence base does its
  work, and where the clipping rule of \Cref{sec:cooling} constrains what may
  be claimed.

  \item[Layer 3 --- Impact and feasibility: \emph{what does it cost, and what
  else does it do?}] Cooling-energy savings are derived from the Layer 2
  temperature estimate and converted to avoided greenhouse-gas emissions;
  capital cost, where supplied, yields payback and a derived cost-feasibility
  score; and co-benefit and equity scores capture what a purely thermal
  assessment would miss.
\end{description}

The three layers are then aggregated into a single NbS Cooling Opportunity
Score, accompanied by per-block confidence ratings, any suitability or
evidence flags, and a deterministically composed recommendation.

\begin{figure}[htbp]
\centering
\resizebox{\textwidth}{!}{% Three-layer conceptual framework.
% Layer 1 Baseline -> Layer 2 NbS performance -> Layer 3 impact & feasibility,
% collected on a right-hand bus into the NbS Cooling Opportunity Score.
% Explicit coordinates (mm) are used so the routing cannot drift.
\begin{tikzpicture}[
  x=1mm, y=1mm,
  font=\small,
  box/.style={draw=rulegrey, line width=0.6pt, rounded corners=1.5pt,
    align=center, inner sep=4pt, fill=white},
  inbox/.style={box, text width=30mm},
  outbox/.style={box, text width=34mm, draw=criterragreen, line width=0.9pt},
  finalbox/.style={box, text width=72mm, draw=criterragreen, line width=1.2pt,
    fill=criterrapaper},
  layerlabel/.style={font=\small\bfseries\color{criterraink}, align=left,
    anchor=north west},
  layerq/.style={font=\footnotesize\itshape\color{criterraink}, align=left,
    anchor=north west},
  flow/.style={-{Latex[length=2mm,width=1.6mm]}, draw=criterraink, line width=0.6pt},
  link/.style={draw=criterraink, line width=0.6pt},
  layerbg/.style={draw=rulegrey, dash pattern=on 2pt off 2pt, line width=0.5pt,
    rounded corners=2pt},
]

\def\junc{28}   % x of the merge junction between inputs and layer output
\def\outx{60}   % x of the layer output box centre
\def\bus{97}    % x of the right-hand collection bus

%% ================= Layer 1 =================
\node[inbox] (l1a) at (0,0)   {Heat Exposure\\Score};
\node[inbox] (l1b) at (0,-14) {Vulnerability\\Score};
\node[outbox] (hpi) at (\outx,-7) {\textbf{Heat Priority}\\\textbf{Index} (0--100)};

\draw[link] (l1a.east) -- (\junc,0);
\draw[link] (l1b.east) -- (\junc,-14);
\draw[link] (\junc,0) -- (\junc,-14);
\draw[flow] (\junc,-7) -- (hpi.west);

%% ================= Layer 2 =================
\node[inbox] (l2a) at (0,-42) {Typology library\\value};
\node[inbox] (l2b) at (0,-56) {Site adjustment\\factor};
\node[outbox] (cps) at (\outx,-49)
  {\textbf{Cooling Potential}\\\textbf{Score} (0--100)\\[1pt]
   \footnotesize{}+ $\Delta T$ range (\si{\degreeCelsius})};

\draw[link] (l2a.east) -- (\junc,-42);
\draw[link] (l2b.east) -- (\junc,-56);
\draw[link] (\junc,-42) -- (\junc,-56);
\draw[flow] (\junc,-49) -- (cps.west);

%% ================= Layer 3 =================
\node[inbox] (l3a) at (0,-86)  {Energy savings\\$\rightarrow$ GHG avoided};
\node[inbox] (l3b) at (0,-100) {Capital cost\\$\rightarrow$ payback};
\node[inbox] (l3c) at (0,-114) {Co-benefits\\and equity};
\node[outbox] (imp) at (\outx,-100)
  {\textbf{Impact and}\\\textbf{feasibility}\\scores};

\draw[link] (l3a.east) -- (\junc,-86);
\draw[link] (l3b.east) -- (\junc,-100);
\draw[link] (l3c.east) -- (\junc,-114);
\draw[link] (\junc,-86) -- (\junc,-114);
\draw[flow] (\junc,-100) -- (imp.west);

%% ================= Collection bus =================
\draw[link] (hpi.east) -- (\bus,-7);
\draw[link] (cps.east) -- (\bus,-49);
\draw[link] (imp.east) -- (\bus,-100);
\draw[link] (\bus,-7) -- (\bus,-134);

\node[finalbox] (final) at (36,-134)
  {\textbf{NbS Cooling Opportunity Score} (0--100)\\[1pt]
   \footnotesize with per-block confidence, flags, and recommendation};

\draw[flow] (\bus,-134) -- (final.east);

%% ================= Layer annotations =================
\node[layerlabel] (lab1) at (-62,4)
  {Layer 1\\[1pt]\normalfont\footnotesize Baseline assessment};
\node[layerq] at (-62,-9) {Does this site\\deserve attention?};

\node[layerlabel] (lab2) at (-62,-38)
  {Layer 2\\[1pt]\normalfont\footnotesize NbS performance};
\node[layerq] at (-62,-51) {What can this\\intervention deliver\\on this site?};

\node[layerlabel] (lab3) at (-62,-82)
  {Layer 3\\[1pt]\normalfont\footnotesize Impact and feasibility};
\node[layerq] at (-62,-95) {What does it cost,\\and what else\\does it do?};

\begin{scope}[on background layer]
  \node[layerbg, fit=(lab1)(l1a)(l1b)(hpi), inner sep=3mm] {};
  \node[layerbg, fit=(lab2)(l2a)(l2b)(cps), inner sep=3mm] {};
  \node[layerbg, fit=(lab3)(l3a)(l3c)(imp), inner sep=3mm] {};
\end{scope}

\end{tikzpicture}
}
\caption[The three-layer conceptual framework]{The three-layer conceptual
framework. Layer 1 characterises the site independently of any intervention;
Layer 2 evaluates one proposed intervention against that site; Layer 3 adds
economic, climate, and social dimensions. The layers are aggregated by the
weighted sum of \Cref{sec:aggregation}. Because Layer 1 does not depend on the
intervention, two options assessed for the same site differ only from Layer 2
onward.}
\label{fig:framework}
\end{figure}

\subsection{Why this decomposition}

The separation is not merely presentational; it has three consequences that
matter for how outputs should be read.

\textbf{It separates the place from the project.} A high Heat Priority Index
with a low Cooling Potential Score identifies a site that badly needs
intervention but for which the proposed intervention is poorly matched --- a
result that a single blended score would conceal. This pattern is common and
decision-relevant: it argues for a different intervention, not for abandoning
the site.

\textbf{It keeps the evidence base localised.} Layer 2 is where published
cooling evidence enters, and it enters in one place: the typology library of
\Cref{sec:typologies}. Revising the evidence base means revising that library
and its adjustment rules, not tracing consequences through the whole
methodology. This is what makes the version-controlled configuration described
in \Cref{sec:implementation} tractable.

\textbf{It makes the confidence structure natural.} A site typically has good
physical data and poor economic data, or the reverse. Because the layers are
distinct, confidence can be computed and reported per block
(\Cref{sec:confidence}) rather than collapsed into a single figure that would
be dominated by whichever block happened to be worst.

\subsection{From inputs to outputs}

\Cref{fig:dataflow} shows the same structure as an implementation data flow.
Three boundaries in that diagram carry the design commitments of
\Cref{sec:commitments} into the software.

First, methodology values reach the engine only as configuration. The
scoring code contains no typology value, no weight, and no threshold; all of
these live in \texttt{config/*.yaml} and are validated on load. Changing the
methodology never requires changing code, and the methodology version travels
with the result.

Second, the engine is a pure function of its inputs and that configuration. It
performs no input or output, makes no network call, holds no global state, and
contains no stochastic component. Determinism is therefore a structural
property rather than a testing aspiration.

Third, every output is a range or a categorical judgment accompanied by its
confidence, and every default the engine applied is itemised in the result. A
reader of a report can see exactly which numbers came from the user and which
from the tool.

\begin{figure}[htbp]
\centering
\resizebox{\textwidth}{!}{% Input -> engine -> output data flow, showing the configuration-as-data
% boundary and the deterministic pure-function contract of the engine.
% Explicit coordinates (mm) keep the routing stable.
\begin{tikzpicture}[
  x=1mm, y=1mm,
  font=\footnotesize,
  box/.style={draw=rulegrey, line width=0.6pt, rounded corners=1.5pt,
    align=center, inner sep=4pt, fill=white},
  stage/.style={box, text width=38mm},
  cfg/.style={box, text width=38mm, draw=criterraink,
    dash pattern=on 2pt off 1.5pt},
  engine/.style={box, text width=40mm, draw=criterragreen, line width=1pt},
  resultbox/.style={box, text width=36mm, draw=criterragreen, line width=0.8pt},
  grouplabel/.style={font=\footnotesize\bfseries\color{criterraink}},
  flow/.style={-{Latex[length=2mm,width=1.6mm]}, draw=criterraink, line width=0.6pt},
  link/.style={draw=criterraink, line width=0.6pt},
  cfgflow/.style={-{Latex[length=2mm,width=1.6mm]}, draw=criterraink,
    line width=0.6pt, dash pattern=on 2pt off 1.5pt},
  groupbg/.style={draw=rulegrey, dash pattern=on 2pt off 2pt, line width=0.5pt,
    rounded corners=2pt},
]

\def\junc{27}    % merge junction for the six input groups
\def\valx{56}    % validation stage centre
\def\engx{104}   % engine centre
\def\outx{158}   % results column centre
\def\fan{132}    % fan-out point to the results column

%% ------------- Inputs (six groups) -------------
\node[stage] (i1) at (0,0)   {1.\ Project information};
\node[stage] (i2) at (0,-13) {2.\ Site characteristics};
\node[stage] (i3) at (0,-26) {3.\ Climate and heat exposure};
\node[stage] (i4) at (0,-39) {4.\ Vulnerability and equity};
\node[stage] (i5) at (0,-52) {5.\ Proposed NbS intervention};
\node[stage] (i6) at (0,-65) {6.\ Cost and energy\\\emph{(optional)}};
\node[grouplabel] (ilab) at (0,10) {Structured inputs};

\foreach \y in {0,-13,-26,-39,-52,-65}{\draw[link] (21,\y) -- (\junc,\y);}
\draw[link] (\junc,0) -- (\junc,-65);

%% ------------- Validation -------------
\node[stage, text width=28mm] (val) at (\valx,-32.5)
  {Schema validation\\[1pt]\emph{errors block,}\\\emph{warnings do not}};
\draw[flow] (\junc,-32.5) -- (val.west);

%% ------------- Engine -------------
\node[engine] (eng) at (\engx,-32.5)
  {\textbf{Scoring engine}\\[1pt]
   pure $\cdot$ deterministic $\cdot$ no I/O\\[1pt]
   \texttt{run\_assessment()}};
\draw[flow] (val.east) -- (eng.west);

%% ------------- Configuration -------------
\node[cfg] (cfg1) at (\engx,4)
  {\texttt{config/*.yaml}\\[1pt]
   typologies $\cdot$ weights\\
   adjustment factors\\
   input mapping $\cdot$ energy model};
\node[grouplabel] at (\engx,17) {Methodology as data};
\draw[cfgflow] (cfg1.south) -- (eng.north)
  node[midway, right=1.5mm, font=\scriptsize\itshape, align=left]
  {version\\\methver};

%% ------------- Outputs -------------
\node[resultbox] (o1) at (\outx,6)   {Heat Priority Index};
\node[resultbox] (o2) at (\outx,-7)  {Cooling Potential Score + $\Delta T$ range};
\node[resultbox] (o3) at (\outx,-20) {Energy savings and GHG avoided};
\node[resultbox] (o4) at (\outx,-33) {Costs and payback};
\node[resultbox] (o5) at (\outx,-46) {Co-benefits and equity};
\node[resultbox] (o6) at (\outx,-59) {\textbf{Opportunity Score}};
\node[resultbox] (o7) at (\outx,-72) {Per-block confidence, flags, assumptions};
\node[grouplabel] (olab) at (\outx,19) {Results (ranges, never point estimates)};

\draw[link] (eng.east) -- (\fan,-32.5);
\draw[link] (\fan,6) -- (\fan,-72);
\foreach \y in {6,-7,-20,-33,-46,-59,-72}{\draw[flow] (\fan,\y) -- (139,\y);}

\begin{scope}[on background layer]
  \node[groupbg, fit=(ilab)(i1)(i6), inner sep=3mm] {};
  \node[groupbg, fit=(olab)(o1)(o7), inner sep=3mm] {};
\end{scope}

%% ------------- Determinism note -------------
\node[font=\scriptsize\itshape\color{criterraink}, align=center,
      text width=150mm] at (72,-92)
  {Identical inputs and an identical methodology version always produce
   identical outputs. Every result records the methodology version that
   produced it, and itemises every default the engine applied.};

\end{tikzpicture}
}
\caption[Input, engine, and output data flow]{Data flow from the six input
groups through validation and the scoring engine to the result. Methodology
values enter only as versioned configuration (dashed path); the engine is a
pure, deterministic function of inputs and configuration; and every output is
reported as a range or a categorical judgment with its own confidence rating.
Validation distinguishes errors, which block an assessment, from warnings,
which do not.}
\label{fig:dataflow}
\end{figure}
